\documentclass{article}
\usepackage{geometry}
\usepackage{graphicx} % Required for inserting images
\usepackage{amsmath, amsthm, amssymb}
\usepackage{parskip}
\newgeometry{vmargin={15mm}, hmargin={24mm,34mm}}
\theoremstyle{definition} 
\newtheorem{definition}{Definition}

\newtheorem{theorem}{Theorem}[section]
\newtheorem{lemma}[theorem]{Lemma}
\newtheorem{corollary}{Corollary}[theorem]


\title{MATH110BH Homework 4}
\date{February 2024}
\author{Boran Erol}

\begin{document}

\maketitle

\section{Problem 1}

\begin{lemma}\label{cyclic_module_iso_to_r_mod_annihiliator}
    Let $ M $ be a cyclic module generated by $ x $. Then, 

    \[ M \cong R/\text{Ann}_{R}(x) \]
\end{lemma}
\begin{proof}
    Consider the surjective homomorphism $ \pi $ from $ R $ to $ M $ given by $ r \mapsto rx $.

    By the first isomorphism theorem for modules, $ R/ \ker(\pi) \cong M $.
\end{proof}


\section{Problem 4}

\begin{lemma}
    Let $ F $ be a free $ R $-module with basis $ (f_{i})_{i \in I} $.
    Let $ J $ be a two-sided ideal of $ R $.
    Then, $ (f_{i} + JF)_{i \in I} $ is a basis for the $ R/F $-module $ F/JF $. 
\end{lemma}
\begin{proof}
    Let $ X = (f_{i})_{i \in I} $.
    Let $ \bar{X} = (f_{i} + J)_{i \in I} $.

    If $ F $ has basis $ (f_{i})_{i \in I} $, we have an $ R $-module isomorphism 
    $ R^{(X)} \equiv F $.

    By Problem 3, modding out by $ JF $ is right exact, so the map $ R^{(X)} / J (R^{(X)}) \to F/JF $.
    Notice that $ R^{(X)} / J (R^{(X)}) \equiv (R/J)^{(\bar{X})} $. Thus, $ \bar{X} $ generates $ F/JF $.

    We'll now show linear independence.
    Let $ \bar{x}_{1}, \ldots, \bar{x}_{n} \in \bar{X} $ and $ r_{1} + J,  \ldots, r_{n} + J \in R/J $ such that

    \[ sum_{i = 1}^{n} (r_{i} + J)(\bar{x}_{i}) \in JF \]

    Then, we have $ a_{1}, \ldots a_{m} \in J $ and $ y_{1}, \ldots, y_{m} \in F $ such that

    \[ \sum_{i = 1}^{n} r_{i}x_{i} = \sum_{i = 1}^{m} a_{i}y_{i} \]

    Moreover, we have $ y_{ij} \in X $ such that $ y_{i} = \sum_{j = 1}^{n_{i}} b_{ij} f_{ij} $.

    But then, since $ X $ is a basis, we can equate all of these coefficients. But this implies that $ r_{i} $
    is a linear combination of the $ a_{j} $ for every $ j $, and thus they are all in $ J $, i.e. $ r_{i} + J = 0 $
    in $ R/J $. This is what we wanted to show.
\end{proof}

\section{Problem 7}

(The first part of the argument is in Notability.)

Let $ I = \{ n \in \mathbb{Z} : ny \in Y^{\prime} \} $. $ I $ is an ideal of $ \mathbb{Z} $,
which is a PID, and is thus generated by a unique $ n > 0 $. This is the smallest $ n $ such that
$ yn \in Y^{\prime} $.

Notice that $ I $ can't be unit a ideal since that would imply $ y \in Y^{\prime} $, which is a contradiction.

If $ I $ is the zero ideal, then $ \langle Y^{\prime}, y \rangle = Y^{\prime} \oplus \langle y \rangle $.
But then, we can use the zero morphism from $ \langle y \rangle $ to induce a map from the direct sum using
the universal property.

Otherwise, we have that $ ny = b $ for some $ b \in Y^{\prime} $.

Since $ A $ is divisible, there exists some $ a \in A $ such that $ na = f(b) $.
We'd like to define $ \bar{f} :  \langle Y^{\prime}, y \rangle \to A $ by mapping
$ t y^{\prime} + sy $ to $ t f(y^{\prime}) + s f(a) $ where $ t,s \in \mathbb{Z} $.

We first need to show that this is well-defined.

Let $ t_{1} y^{\prime}_{1} + s_{1}y = t_{2} y^{\prime}_{2} + s_{2}y $.

Notice that $ (s_{1} - s_{2})a \in Y^{\prime} $ so $ s_{1} - s_{2} $ is divisible by $ n $.

Then, $ f(t_{1} y^{\prime}_{1} + s_{1}y) = t_{1} f(y^{\prime}_{1}) + s_{1}a $.
Similarly, $ f(t_{2} y^{\prime}_{2} + s_{2}y) = t_{2} f(y^{\prime}_{2}) + s_{2}a $.

Then, $ f(t_{1} y^{\prime}_{1} + t_{2} y^{\prime}_{2}) + s_{1}a - s_{2}a = f(t_{1} y^{\prime}_{1} + t_{2} y^{\prime}_{2}) + nma = f(t_{1} y^{\prime}_{1} + t_{2} y^{\prime}_{2}) + f(mb) $.

But now $ f $ maps zero to zero, and thus we're done.

This is clearly a homomorphism, so we've extended $ f $, which is a contradiction. We thus conclude the proof.

\section{Problem 8}

\begin{lemma}
    The quotient of a injective object is a injective object in any category.
\end{lemma}
\begin{proof}
    First of all, notice that this is the dual statement to the subgroup of a free group being a free group,
    which we know to be true. Another way to see this is to use the extension definition of the injective
    element and lift the surjection to the previous element to extend it.
\end{proof}

Also notice that the product of injective groups is injective since we can extend them one by one using the projection
maps and then use the universal property to put them back together.

Let $ A $ be an Abelian group. Then, $ A \equiv F/G $ for some free Abelian group $ F $.

Notice that $ F = \oplus_{i \in I} \mathbb{Z} $ which is a subgroup of $ \oplus_{i \in I} \mathbb{Q} $, which is injective.

Then, $ A \equiv F/G $ injects into a quotient of $ \oplus_{i \in I} \mathbb{Q} $, which is injective by the lemma above.
We thus conclude the proof.

\end{document}
